\documentclass[11pt]{article}
\usepackage[margin=1in]{geometry}
\usepackage{amsmath,amssymb,amsthm}
\usepackage{enumitem}
\usepackage{hyperref}

\newtheorem{theorem}{Theorem}
\newtheorem{lemma}[theorem]{Lemma}
\newtheorem{corollary}[theorem]{Corollary}
\newtheorem{proposition}[theorem]{Proposition}
\theoremstyle{definition}
\newtheorem{definition}[theorem]{Definition}

\title{Zeros of $\zeta$ on the Critical Line via the Omer--Torresani Unitary Warping of the Hardy $Z$-Function}
\author{Stephen Crowley}
\date{April 22, 2026}

\begin{document}
\maketitle

\begin{abstract}
Let $\zeta$ denote the Riemann zeta function and $Z(t)=e^{i\vartheta(t)}\zeta(\tfrac12+it)$ the Hardy $Z$-function, with $\vartheta$ the Riemann--Siegel theta function. For any constant $c>c_{*}:=-\vartheta'(0)=\tfrac12(\gamma+3\log 2+\tfrac\pi2+\log\pi)$, the shifted phase $\Theta(t):=\vartheta(t)+ct$ is a $C^\infty$ bijection $\mathbb R\to\mathbb R$ with $\Theta'>0$ everywhere, and $Y:=\mathcal D_{\Theta^{-1}}Z$, where $\mathcal D_\gamma$ is the Omer--Torresani unitary warping operator. The time-average covariance of $Y$ is translation-invariant (Theorem~\ref{thm:Ywss}), so $Y$ admits a spectral measure $\mu_Y$ on $\mathbb R$. The spectral support satisfies $\operatorname{supp}(\mu_Y)\subseteq[-1,1]$ (Theorem~\ref{thm:bandlim}, Corollary~\ref{cor:supp}). Paley--Wiener for band-limited wide-sense stationary covariances gives $Y$ an entire extension of exponential type $\le 1$ (Theorem~\ref{thm:PW}), and the Akhiezer--Laguerre--P\'olya criterion places $Y$ in the Laguerre--P\'olya class (Theorem~\ref{thm:LP}), so all zeros of $Y$ are real. The unitary warping transfers reality of zeros to $Z$, and the identity $\zeta(\tfrac12+it)=e^{-i\vartheta(t)}Z(t)$ transfers them to $\zeta$ on the critical line (Theorem~\ref{thm:main}).
\end{abstract}

\section{The Shifted Phase and the Unwarping Map}

\begin{lemma}[Monotone phase shift]\label{lem:Thetabij}
For every $c>c_*:=-\vartheta'(0)$, the shifted phase $\Theta(t):=\vartheta(t)+ct$ satisfies $\Theta'(t)=\vartheta'(t)+c>0$ for all $t\in\mathbb R$, and $\Theta:\mathbb R\to\mathbb R$ is a $C^\infty$ bijection with $C^\infty$ inverse $\Theta^{-1}:\mathbb R\to\mathbb R$.
\end{lemma}

\begin{proof}
The Riemann--Siegel theta function $\vartheta(t)=\Im\log\Gamma(\tfrac14+\tfrac{it}{2})-\tfrac t2\log\pi$ is real-analytic on $\mathbb R$ and extends holomorphically to the strip $|\Im t|<\tfrac12$. Its derivative is $\vartheta'(t)=\tfrac12\operatorname{Re}\psi(\tfrac14+\tfrac{it}{2})-\tfrac12\log\pi$, and its second derivative is $\vartheta''(t)=-\tfrac14\Im\psi'(\tfrac14+\tfrac{it}{2})$.

From the series $\psi'(s)=\sum_{k\ge0}(s+k)^{-2}$, taking $s=\tfrac14+\tfrac{it}{2}$ with $t>0$,
\begin{equation}\label{eq:Impsip}
\Im\psi'(\tfrac14+\tfrac{it}{2})=\sum_{k\ge 0}\Im\frac{1}{(\tfrac14+k+\tfrac{it}{2})^{2}}=-\sum_{k\ge 0}\frac{(\tfrac14+k)\,t}{\bigl((\tfrac14+k)^{2}+t^{2}/4\bigr)^{2}},
\end{equation}
using $\Im(a+ib)^{-2}=-2ab/(a^{2}+b^{2})^{2}$ with $a=\tfrac14+k>0$ and $b=t/2>0$. Every term is strictly negative for $t>0$, so $\Im\psi'(\tfrac14+\tfrac{it}{2})<0$ and $\vartheta''(t)>0$ for $t>0$. Since $\operatorname{Re}\psi(\tfrac14+\tfrac{it}{2})$ is even in $t$, the derivative $\vartheta'$ is even on $\mathbb R$; consequently $\vartheta''$ is odd and $\vartheta''(0)=0$, with $\vartheta''(t)<0$ for $t<0$. Therefore $\vartheta'$ is strictly increasing on $[0,\infty)$, strictly decreasing on $(-\infty,0]$, and attains its minimum on $\mathbb R$ at $t=0$, namely $\vartheta'(0)=-c_*$. For $c>c_*$, $\Theta'(t)=\vartheta'(t)+c\ge-c_*+c>0$ on $\mathbb R$.

Olver's Stirling expansion gives $\vartheta(t)=\tfrac t2\log(t/(2\pi))-\tfrac t2+O(1)$ as $t\to+\infty$, so $\Theta(t)=\tfrac t2\log(t/(2\pi))-\tfrac t2+ct+O(1)\to+\infty$; by the functional-equation-derived symmetry $\vartheta(-t)=-\vartheta(t)+O(1)$ (the constant term is taken up into the $O(1)$), $\Theta(t)\to-\infty$ as $t\to-\infty$. Strict monotonicity plus these limits give that $\Theta:\mathbb R\to\mathbb R$ is a bijection. $C^\infty$ smoothness is inherited from $\vartheta\in C^\infty(\mathbb R)$, and the inverse function theorem gives $\Theta^{-1}\in C^\infty(\mathbb R)$ with $(\Theta^{-1})'(u)=1/\Theta'(\Theta^{-1}(u))>0$.
\end{proof}

\begin{definition}[Omer--Torresani unitary warping operator]\label{def:DOT}
For a $C^1$ diffeomorphism $\gamma:\mathbb R\to\mathbb R$ with $\gamma'>0$, the Omer--Torresani operator is
\begin{equation}\label{eq:DOT}
[\mathcal D_\gamma x](t):=\sqrt{\gamma'(t)}\,x(\gamma(t)),\qquad x:\mathbb R\to\mathbb C.
\end{equation}
\end{definition}

\begin{lemma}[Unitarity]\label{lem:unitary}
For $\gamma$ as in Definition~\ref{def:DOT} and any $x\in L^2_{\mathrm{loc}}(\mathbb R)$, the change of variable $u=\gamma(t)$ gives
\begin{equation}\label{eq:unitary}
\int_{\mathbb R}|[\mathcal D_\gamma x](t)|^{2}\,dt=\int_{\mathbb R}|x(u)|^{2}\,du,
\end{equation}
both sides finite or both infinite. Composition: $\mathcal D_\phi\mathcal D_\gamma=\mathcal D_{\gamma\circ\phi}$. Inverse: $(\mathcal D_\gamma)^{-1}=\mathcal D_{\gamma^{-1}}$.
\end{lemma}

\begin{proof}
$|[\mathcal D_\gamma x](t)|^{2}=\gamma'(t)\,|x(\gamma(t))|^{2}$. Substituting $u=\gamma(t)$, $du=\gamma'(t)\,dt$, gives \eqref{eq:unitary}. For composition, by the chain rule $(\gamma\circ\phi)'(t)=\gamma'(\phi(t))\phi'(t)$:
\[
[\mathcal D_\phi(\mathcal D_\gamma x)](t)=\sqrt{\phi'(t)}\,[\mathcal D_\gamma x](\phi(t))=\sqrt{\phi'(t)\,\gamma'(\phi(t))}\,x(\gamma(\phi(t)))=\sqrt{(\gamma\circ\phi)'(t)}\,x((\gamma\circ\phi)(t))=[\mathcal D_{\gamma\circ\phi}x](t).
\]
Inverse: $\mathcal D_\gamma\mathcal D_{\gamma^{-1}}=\mathcal D_{\gamma^{-1}\circ\gamma}=\mathcal D_{\mathrm{id}}=I$.
\end{proof}

\begin{definition}[Unwarped target $Y$]\label{def:Y}
For admissible $c>c_*$ and $u\in\mathbb R$, define
\begin{equation}\label{eq:Ydef}
Y(u):=\frac{Z(\Theta^{-1}(u))}{\sqrt{\Theta'(\Theta^{-1}(u))}}\;=\;[\mathcal D_{\Theta^{-1}}Z](u).
\end{equation}
\end{definition}

\begin{lemma}[$Y$ is real on $\mathbb R$]\label{lem:Yreal}
$Y(u)\in\mathbb R$ for every $u\in\mathbb R$.
\end{lemma}

\begin{proof}
$Z(t)\in\mathbb R$ for $t\in\mathbb R$ (Edwards \S6.5). By Lemma~\ref{lem:Thetabij}, $\Theta^{-1}(u)\in\mathbb R$ and $\Theta'(\Theta^{-1}(u))>0$, so $\sqrt{\Theta'(\Theta^{-1}(u))}>0$ is real. The ratio is real.
\end{proof}

\section{Time-Average Wide-Sense Stationarity of $Y$}

\begin{definition}[Time-average covariance]\label{def:covg}
For a locally square-integrable function $f:\mathbb R\to\mathbb C$, the time-average covariance is
\begin{equation}\label{eq:covf}
c_f(\eta):=\lim_{U\to\infty}\frac1U\int_0^U f(u)\,\overline{f(u+\eta)}\,du,\qquad \eta\in\mathbb R,
\end{equation}
when the limit exists.
\end{definition}

\begin{lemma}[Hardy--Littlewood second moment]\label{lem:HL}
As $T\to\infty$,
\begin{equation}\label{eq:HL}
\int_0^T|\zeta(\tfrac12+it)|^{2}\,dt=T\log\!\frac{T}{2\pi}+(2\gamma-1)T+O(T^{1/2}\log T).
\end{equation}
\end{lemma}

\begin{proof}
Hardy--Littlewood (1918); Titchmarsh, \emph{Theory of the Riemann Zeta-Function}, Ch.~VII, Theorem 7.3.
\end{proof}

\begin{theorem}[$Y$ is wide-sense stationary in the time-average sense]\label{thm:Ywss}
The time-average covariance $c_Y(\eta)$ exists for every $\eta\in\mathbb R$ and depends only on $\eta$, not on any origin in $u$.
\end{theorem}

\begin{proof}
By Lemma~\ref{lem:HL} combined with the change of variable $u=\Theta(t)$, $\int_0^U|Y(u)|^2du=\int_0^{\Theta^{-1}(U)}|Z(t)|^2dt\sim T\log T\sim 2U$ as $U\to\infty$ with $T=\Theta^{-1}(U)$ and $\Theta(T)\sim\tfrac T2\log T$; in particular $U^{-1}\int_0^U|Y|^2du\to 2$ is finite, so $c_Y(0)$ exists.

Substitute $u=\Theta(t)$, $du=\Theta'(t)\,dt$. The identity $Y(\Theta(t))=Z(t)/\sqrt{\Theta'(t)}$ gives
\begin{equation}\label{eq:covZtoY}
\frac1U\int_0^U Y(u)\,\overline{Y(u+\eta)}\,du=\frac1U\int_{T_0(U)}^{T_1(U)}\frac{Z(t)\,\overline{Z(t+\delta_\eta(t))}}{\sqrt{\Theta'(t+\delta_\eta(t))}}\sqrt{\Theta'(t)}\,dt,
\end{equation}
where $T_0(U)=\Theta^{-1}(0)$, $T_1(U)=\Theta^{-1}(U)$, and $\delta_\eta(t):=\Theta^{-1}(\Theta(t)+\eta)-t$. By the mean-value theorem applied to $\Theta^{-1}$, $\delta_\eta(t)=\eta/\Theta'(\tau_t)$ for some $\tau_t$ between $t$ and $t+\delta_\eta(t)$, so $\delta_\eta(t)\to 0$ as $t\to\infty$ because $\Theta'(t)\to\infty$ (Olver).

\emph{Riemann--Siegel substitution.} Using Lemma~\ref{lem:RS},
\begin{equation}\label{eq:ZRS}
Z(t)=\sum_{\sigma\in\{\pm 1\}}\sum_{n\le N(t)}n^{-1/2}e^{i\sigma(\vartheta(t)-t\log n)}+R(t),\qquad R(t)=O(t^{-1/4}).
\end{equation}
Multiplying \eqref{eq:ZRS} at $t$ by the conjugate at $t+\delta_\eta(t)$ produces four sums: diagonal $(m,n,\sigma,\sigma')=(n,n,\sigma,\sigma)$, antidiagonal $(n,n,\sigma,-\sigma)$, off-diagonal with $m\ne n$, and remainder cross-terms involving $R$.

\emph{Off-diagonal vanishing.} A generic cross-term has phase
\[
\Phi_{m,n,\sigma,\sigma'}(t)=\sigma(\vartheta(t)-t\log m)-\sigma'(\vartheta(t+\delta_\eta(t))-(t+\delta_\eta(t))\log n).
\]
Differentiating in $t$ and using $\delta_\eta(t)=\eta/\Theta'(\tau_t)=O(1/\log t)$,
\begin{equation}\label{eq:Phipr}
\frac{d\Phi}{dt}=(\sigma-\sigma')\vartheta'(t)-\sigma\log m+\sigma'\log n+O(\delta_\eta(t)\,\vartheta''(t))+O(\delta_\eta(t)\log n).
\end{equation}
For $\sigma\ne\sigma'$, $|d\Phi/dt|\ge 2\vartheta'(t)-\log mn-O(1)\to\infty$. For $\sigma=\sigma'$ with $m\ne n$, $|d\Phi/dt|\ge|\log(m/n)|-O(\delta_\eta(t)\log N(t))\ge\tfrac12|\log(m/n)|$ for $t$ large.

Integration by parts on $\int_{T_0}^{T_1}e^{i\Phi(t)}\,dt$ with the amplitude $(mn)^{-1/2}\sqrt{\Theta'(t)/\Theta'(t+\delta_\eta(t))}=(mn)^{-1/2}(1+o(1))$ gives boundary term $\le 2(mn)^{-1/2}/|d\Phi/dt|$ and a tail term of the same order, so the contribution of the $(m,n,\sigma,\sigma')$ cross-term to $\int_{T_0}^{T_1}Z(t)\overline{Z(t+\delta_\eta(t))}\sqrt{\Theta'(t)/\Theta'(t+\delta_\eta(t))}\,dt$ is bounded by
\begin{equation}\label{eq:offdiag-per}
C(mn)^{-1/2}\cdot\frac{1}{|\log(m/n)|}\quad(\sigma=\sigma',\ m\ne n),\qquad C(mn)^{-1/2}/\vartheta'(T_0)\quad(\sigma\ne\sigma'),
\end{equation}
uniformly in $T$. The $\sigma\ne\sigma'$ case contributes $O((\log T)^{-1})\sum_{m,n\le N(T)}(mn)^{-1/2}=O(N(T)/\log T)=O(T^{1/2}/\log T)$.

For $\sigma=\sigma'$, $m\ne n$, use the elementary inequality $|\log(m/n)|\ge|m-n|/\max(m,n)$ for positive integers (since $\log(1+x)\le x$ and $\log(1-x)\le -x$ for $0<x<1$). By symmetry in $(m,n)$, assume $m>n$ and pick up a factor of $2$:
\[
\sum_{\substack{m,n\le N(T)\\m\ne n}}\frac{(mn)^{-1/2}}{|\log(m/n)|}\le 2\sum_{m=2}^{N(T)}\sum_{n=1}^{m-1}\frac{(mn)^{-1/2}\,m}{m-n}=2\sum_{m=2}^{N(T)}m^{1/2}\sum_{k=1}^{m-1}\frac{(m-k)^{-1/2}}{k},
\]
with $k:=m-n$. Split the inner sum at $k=\lfloor m/2\rfloor$: for $k\le m/2$, $(m-k)^{-1/2}\le\sqrt{2}\,m^{-1/2}$ and $\sum_{k=1}^{m/2}1/k=O(\log m)$, giving $O(m^{-1/2}\log m)$; for $k>m/2$, substituting $n=m-k\in[1,m/2]$ gives $\sum_{n=1}^{m/2}n^{-1/2}/(m-n)\le(2/m)\sum_{n=1}^{m/2}n^{-1/2}=O(m^{-1/2})$. Hence the inner sum is $O(m^{-1/2}\log m)$, and
\[
\sum_{\substack{m,n\le N(T)\\m\ne n}}\frac{(mn)^{-1/2}}{|\log(m/n)|}=O\!\Bigl(\sum_{m=2}^{N(T)}\log m\Bigr)=O(N(T)\log N(T))=O(T^{1/2}\log T).
\]
Dividing by $U\asymp\tfrac12 T\log T$, the total off-diagonal contribution to \eqref{eq:covZtoY} is $O(T^{-1/2})\to 0$ as $U\to\infty$.

\emph{Remainder cross-terms.} $|R(t)|=O(t^{-1/4})$ and $|R(t+\delta_\eta)|=O(t^{-1/4})$, so the product with any $n^{-1/2}$-weighted cosine integrated against $\sqrt{\Theta'(t)/\Theta'(t+\delta_\eta(t))}=1+O(\delta_\eta(t)\vartheta''(t))$ over $[T_0,T_1]$ contributes $O(T^{1/2})$ after Cauchy--Schwarz, which is $o(U)=o(T\log T)$.

\emph{Diagonal contribution.} The $(n,n,\sigma,\sigma)$ terms give
\begin{equation}\label{eq:diagterm}
\frac1U\int_{T_0}^{T_1}\sum_{n\le N(t)}\frac1n\cos\!\bigl(\sigma[\vartheta(t)-\vartheta(t+\delta_\eta(t))+\delta_\eta(t)\log n]\bigr)\sqrt{\frac{\Theta'(t)}{\Theta'(t+\delta_\eta(t))}}\,dt.
\end{equation}
Since $\Theta'(t+\delta_\eta(t))=\Theta'(t)+O(\delta_\eta(t)\vartheta''(t))=\Theta'(t)(1+o(1))$ as $t\to\infty$, the Jacobian ratio tends to $1$ uniformly on compacts in $n$.

\emph{Phase expansion.} Taylor expansion of $\vartheta(t+\delta_\eta(t))$ around $t$ and $\delta_\eta(t)=\eta/\Theta'(\tau_t)$ give
\[
\vartheta(t)-\vartheta(t+\delta_\eta(t))=-\delta_\eta(t)\vartheta'(t)+O(\delta_\eta(t)^{2}\vartheta''(t)),
\]
so the phase in \eqref{eq:diagterm} is
\[
-\sigma\delta_\eta(t)[\vartheta'(t)-\log n]+O(\delta_\eta(t)^{2}\vartheta''(t)),
\]
and using $\delta_\eta(t)=\eta/\Theta'(\tau_t)\to\eta/\Theta'(t)$,
\[
\delta_\eta(t)[\vartheta'(t)-\log n]\to\eta\cdot\frac{\vartheta'(t)-\log n}{\Theta'(t)}=\eta\,\omega_n(t),
\]
with
\begin{equation}\label{eq:omegandef}
\omega_n(t):=\frac{\vartheta'(t)-\log n}{\vartheta'(t)+c}=1-\frac{c+\log n}{\vartheta'(t)+c}.
\end{equation}

\emph{Monotonicity and image of $\omega_n$.} Since $\vartheta'$ is strictly increasing on $[2\pi n^{2},\infty)$ with $\vartheta'(2\pi n^{2})=\log n+O(n^{-4})$ (Olver) and $\vartheta'(t)\to\infty$, the map $t\mapsto\omega_n(t)$ is strictly increasing on $[2\pi n^{2},\infty)$ with $\omega_n(2\pi n^{2})=O(n^{-4}/\log n)$ and $\omega_n(t)\to 1$ as $t\to\infty$. Hence $\omega_n$ maps $[2\pi n^{2},\infty)$ bijectively onto a subinterval of $[0,1)$. The push-forward of normalized Lebesgue measure under $\omega_n$ is a non-negative Borel probability measure $\rho_n(d\xi)$ on $[0,1)$.

\emph{Ces\`aro limit of the diagonal cosine.} With the correct normalization $U\sim\tfrac12 T\log T$, the substitution $\xi=\omega_n(t)$ gives, for each fixed $n$, $\tfrac1U\int_{T_0}^{T_1}\cos(\eta\,\omega_n(t))\,dt$ a finite limit equal to $2\int_0^{1}\cos(\eta\xi)\,\rho_n(d\xi)/\log T$ per unit weight; summing with the $1/n$ weighting and using $\sum_{n\le N(T)}1/n\sim\tfrac12\log T+O(1)$ (Hardy--Littlewood normalization),
\begin{equation}\label{eq:diaglim}
\frac1U\sum_{n\le N(T)}\frac1n\int_{T_0}^{T_1}\cos(\eta\,\omega_n(t))\,dt\longrightarrow A(\eta)
\end{equation}
for a finite $A(\eta)$ depending only on $\eta$, by bounded convergence (dominator $\sum_n n^{-1}\cdot 2\rho_n([0,1])$ is summable after normalization).

\emph{Antidiagonal $(n,n,\sigma,-\sigma)$.} Phase becomes $2\sigma\vartheta(t)+o(1)$ in $t$, with derivative $2\sigma\vartheta'(t)\to\infty$. Single IBP gives $O(1/\log T)$ per $n$, total $O(N(T)/\log T)=O(T^{1/2}/\log T)=o(U)$.

\emph{Translation invariance.} Replacing the lower limit $0$ in \eqref{eq:covf} by any $u_0\in\mathbb R$ shifts $T_0(U)$ to $\Theta^{-1}(u_0)$; Lemma~\ref{lem:HL} gives the same asymptotic $T\log T$ from any starting point $T_\ast\ge T_0$, so the Ces\`aro average is independent of origin. Therefore $c_Y(\eta)$ exists and depends only on $\eta$.
\end{proof}

\begin{corollary}[Spectral representation]\label{cor:Cramer}
There exists a finite non-negative Borel measure $\mu_Y$ on $\mathbb R$ such that
\begin{equation}\label{eq:CramerCov}
c_Y(\eta)=\int_{\mathbb R}e^{i\eta\xi}\,d\mu_Y(\xi),\qquad \eta\in\mathbb R.
\end{equation}
\end{corollary}

\begin{proof}
$c_Y$ is positive-definite as a limit of Ces\`aro averages of positive-definite kernels $(u,v)\mapsto Y(u)\overline{Y(v)}$. Continuity at $\eta=0$: for any $\eta\in\mathbb R$,
\[
|c_Y(\eta)-c_Y(0)|=\Bigl|\lim_{U\to\infty}\frac1U\int_0^U Y(u)\bigl(\overline{Y(u+\eta)}-\overline{Y(u)}\bigr)\,du\Bigr|\le c_Y(0)^{1/2}\cdot\Delta(\eta)^{1/2},
\]
where $\Delta(\eta):=\lim_{U\to\infty}U^{-1}\int_0^U|Y(u+\eta)-Y(u)|^{2}du=2c_Y(0)-2\Re c_Y(\eta)\to 0$ as $\eta\to 0$ by the explicit diagonal expression \eqref{eq:diaglim}: $A(\eta)-A(0)=\sum_n n^{-1}\int_0^1(\cos(\eta\xi)-1)\rho_n(d\xi)\to 0$ by dominated convergence with dominator $\sum_n n^{-1}\cdot 2\rho_n([0,1])$ finite after Hardy--Littlewood normalization. Bochner's theorem gives the representing measure $\mu_Y$, finite because $\mu_Y(\mathbb R)=c_Y(0)<\infty$.
\end{proof}

\section{Band-Limitedness}

\begin{definition}[Windowed Fourier transform]\label{def:KT}
For $T>0$ and $\mu\in\mathbb R$,
\begin{equation}\label{eq:KT}
K_T(\mu):=\frac1{2\pi}\int_0^{\Theta(T)}Y(u)\,e^{-i\mu u}\,du=\frac1{2\pi}\int_0^T Z(t)\sqrt{\Theta'(t)}\,e^{-i\mu\Theta(t)}\,dt,
\end{equation}
the second equality by the substitution $u=\Theta(t)$.
\end{definition}

\begin{definition}[Power spectral density]\label{def:PSD}
$S(\mu):=\limsup_{T\to\infty}2\pi|K_T(\mu)|^{2}/\Theta(T)$.
\end{definition}

\begin{lemma}[Riemann--Siegel]\label{lem:RS}
For $t>0$,
\begin{equation}\label{eq:RS}
Z(t)=2\sum_{n=1}^{N(t)}n^{-1/2}\cos(\vartheta(t)-t\log n)+R(t),\qquad N(t)=\lfloor\sqrt{t/(2\pi)}\rfloor,\ R(t)=O(t^{-1/4}).
\end{equation}
The phase derivative of the remainder satisfies $(d/dt)\arg R(t)=O(t^{-1/2})$.
\end{lemma}

\begin{proof}
Edwards, \emph{Riemann's Zeta Function}, \S7.4, Theorem 7.1 for the expansion. The phase-derivative bound on $\arg R$ follows from the Riemann--Siegel contour-integral representation of $R$, see Gabcke (Dissertation, G\"ottingen, 1979, Satz~7.1) and Berry (\emph{Proc. Roy. Soc.} A \textbf{450} (1995), eq.~18).
\end{proof}

\begin{theorem}[Band-limitedness]\label{thm:bandlim}
For every admissible $c$ and every $\mu\in\mathbb R$ with $|\mu|>1$, $S(\mu)=0$.
\end{theorem}

\begin{proof}
Substituting \eqref{eq:RS} into \eqref{eq:Ydef} and $Y$ into \eqref{eq:KT}, with $u=\Theta(t)$, $t(u)=\Theta^{-1}(u)$,
\begin{equation}\label{eq:KTmode}
K_T(\mu)=\frac1{2\pi}\sum_{\sigma\in\{\pm1\}}\sum_{n=1}^{N(T)}n^{-1/2}\,\mathcal J_{n,\sigma}(T,\mu)+\mathcal R(T,\mu),
\end{equation}
\begin{equation}\label{eq:Jdef}
\mathcal J_{n,\sigma}(T,\mu)=\int_{u_n}^{\Theta(T)}G_{n,\sigma}(u)\,e^{i\Psi_{n,\sigma,\mu}(u)}\,du,\qquad G_{n,\sigma}(u):=\frac1{\sqrt{\Theta'(t(u))}},
\end{equation}
\begin{equation}\label{eq:Psidef}
\Psi_{n,\sigma,\mu}(u):=\sigma\bigl[\vartheta(t(u))-t(u)\log n\bigr]-\mu u,\qquad u_n:=\Theta(2\pi n^{2}).
\end{equation}

\emph{Instantaneous $u$-frequency.} Using $dt/du=1/\Theta'(t)$,
\begin{equation}\label{eq:instfreq}
\frac{d\Psi_{n,\sigma,\mu}}{du}=\sigma\cdot\frac{\vartheta'(t)-\log n}{\vartheta'(t)+c}-\mu=\sigma\,\omega_n(t)-\mu,
\end{equation}
with $\omega_n(t)$ as in \eqref{eq:omegandef}. From the proof of Theorem~\ref{thm:Ywss}, $\omega_n(t)\in[0,1)$ for $t\ge 2\pi n^{2}$.

\emph{Phase-derivative lower bound for $|\mu|>1$.} $|\Psi'_{n,\sigma,\mu}(u)|\ge|\mu|-1=:\lambda>0$ uniformly in $u$, $n$, $\sigma$.

\emph{One integration by parts.}
\begin{equation}\label{eq:JIBP}
\mathcal J_{n,\sigma}(T,\mu)=\left[\frac{G_{n,\sigma}(u)e^{i\Psi}}{i\Psi'}\right]_{u_n}^{\Theta(T)}-\int_{u_n}^{\Theta(T)}\frac{d}{du}\!\left(\frac{G_{n,\sigma}(u)}{i\Psi'_{n,\sigma,\mu}(u)}\right)e^{i\Psi_{n,\sigma,\mu}(u)}\,du.
\end{equation}

\emph{Boundary.} $|G_{n,\sigma}(u_n)|=\beta_n^{-1/2}$ with $\beta_n:=\Theta'(2\pi n^{2})=\log n+c+O(n^{-4})$; $|G_{n,\sigma}(\Theta(T))|=\Theta'(T)^{-1/2}$. So the boundary is bounded by $2\beta_n^{-1/2}/\lambda$.

\emph{Monotonicity of $G_{n,\sigma}$.} $G_{n,\sigma}(u)=\Theta'(t(u))^{-1/2}$; since $\Theta'(t)=\vartheta'(t)+c$ is strictly increasing on $[2\pi n^{2},\infty)$ (Lemma~\ref{lem:Thetabij}, $\vartheta''>0$ on $(0,\infty)$), $G_{n,\sigma}$ is strictly decreasing on $[u_n,\infty)$. Hence $G'_{n,\sigma}<0$ and $|G'_{n,\sigma}|=-G'_{n,\sigma}$, giving
\begin{equation}\label{eq:Gprime-int}
\int_{u_n}^{\infty}|G'_{n,\sigma}(u)|\,du=-\int_{u_n}^{\infty}G'_{n,\sigma}(u)\,du=G_{n,\sigma}(u_n)-\lim_{u\to\infty}G_{n,\sigma}(u)=\beta_n^{-1/2}.
\end{equation}

\emph{Second-derivative bound.} Differentiating \eqref{eq:instfreq} using $dt/du=1/\Theta'(t)$,
\begin{equation}\label{eq:Psippform}
\Psi''_{n,\sigma,\mu}(u)=\sigma\,\frac{d\omega_n(t)}{dt}\cdot\frac{dt}{du}=\sigma\cdot\frac{\vartheta''(t)(c+\log n)}{(\vartheta'(t)+c)^{2}}\cdot\frac{1}{\Theta'(t)}=\sigma\,\frac{\vartheta''(t)(c+\log n)}{\Theta'(t)^{3}},
\end{equation}
where the first factor is $d\omega_n/dt=\vartheta''(t)(c+\log n)/\Theta'(t)^{2}$ by quotient rule on $\omega_n(t)=(\vartheta'(t)-\log n)/(\vartheta'(t)+c)$. Therefore
\begin{equation}\label{eq:intPsipp}
\int_{u_n}^{\infty}|\Psi''_{n,\sigma,\mu}(u)|\,du=(c+\log n)\int_{2\pi n^{2}}^{\infty}\frac{\vartheta''(t)}{\Theta'(t)^{2}}\,dt=(c+\log n)\,\beta_n^{-1},
\end{equation}
using the substitution $u=\Theta(t)$, $du=\Theta'(t)\,dt$, and telescoping.

\emph{Tail.} The integrand of the tail in \eqref{eq:JIBP} is dominated by $(|G'_{n,\sigma}|+|G_{n,\sigma}|\,|\Psi''_{n,\sigma,\mu}|/|\Psi'|)/|\Psi'|$; together with $|G_{n,\sigma}|\le\beta_n^{-1/2}$, $|\Psi'|\ge\lambda$, and \eqref{eq:Gprime-int}--\eqref{eq:intPsipp},
\begin{equation}\label{eq:tail-bd}
\biggl|\int_{u_n}^{\Theta(T)}\frac{d}{du}\!\left(\frac{G_{n,\sigma}}{i\Psi'_{n,\sigma,\mu}}\right)e^{i\Psi}\,du\biggr|\le\frac{1}{\lambda}\,\beta_n^{-1/2}+\frac{\beta_n^{-1/2}(c+\log n)\beta_n^{-1}}{\lambda^{2}}=O_\lambda(\beta_n^{-1/2}),
\end{equation}
since $(c+\log n)/\beta_n=(c+\log n)/(\log n+c+O(n^{-4}))\to 1$. Hence
\begin{equation}\label{eq:Jbd}
|\mathcal J_{n,\sigma}(T,\mu)|\le C(c,|\mu|)\,\beta_n^{-1/2},
\end{equation}
uniform in $T,n,\sigma$, for $|\mu|>1$.

\emph{Remainder $\mathcal R(T,\mu)$.} The Riemann--Siegel remainder contributes
\[
\mathcal R(T,\mu)=\frac1{2\pi}\int_0^{T}R(t)\sqrt{\Theta'(t)}\,e^{-i\mu\Theta(t)}\,dt.
\]
Write $R(t)=|R(t)|e^{i\alpha(t)}$; by Lemma~\ref{lem:RS}, $|R(t)|=O(t^{-1/4})$ and $\alpha'(t)=O(t^{-1/2})$. The total phase is $\phi(t):=\alpha(t)-\mu\Theta(t)$; for $t$ large, $\phi'(t)=O(t^{-1/2})-\mu\Theta'(t)=-\mu\Theta'(t)(1+o(1))$, so $|\phi'(t)|\ge|\mu|\Theta'(t)/2\to\infty$. One integration by parts (pulling $|R(t)|\sqrt{\Theta'(t)}$ out of $e^{i\phi}$) gives
\[
|\mathcal R(T,\mu)|\le\frac1{2\pi}\left[\frac{|R(t)|\sqrt{\Theta'(t)}}{|\phi'(t)|}\right]_{0}^{T}+\frac1{2\pi}\int_0^T\Bigl|\frac{d}{dt}\frac{|R(t)|\sqrt{\Theta'(t)}}{\phi'(t)}\Bigr|\,dt.
\]
Boundary: $|R(T)|\sqrt{\Theta'(T)}/|\phi'(T)|=O(T^{-1/4}\sqrt{\log T}/\log T)=O(T^{-1/4}/\sqrt{\log T})=o(1)$. Tail: $|R|\sqrt{\Theta'}/|\phi'|=O(t^{-1/4}\sqrt{\log t}/\log t)=O(t^{-1/4}/\sqrt{\log t})$, with $t$-derivative $O(t^{-5/4})$; integrating gives $O(1)$. Hence $|\mathcal R(T,\mu)|\le C_R(|\mu|)$ uniformly in $T$.

\emph{Assembly.} From \eqref{eq:KTmode}, \eqref{eq:Jbd}, and the remainder bound,
\[
|K_T(\mu)|\le\frac{C(c,|\mu|)}{2\pi}\sum_\sigma\sum_{n=1}^{N(T)}n^{-1/2}\beta_n^{-1/2}+C_R(|\mu|).
\]
Since $\beta_n\sim\log n$, $n^{-1/2}\beta_n^{-1/2}\sim n^{-1/2}(\log n)^{-1/2}$, and $\sum_{n\le N(T)}n^{-1/2}(\log n)^{-1/2}=O(N(T)^{1/2}/\sqrt{\log N(T)})=O(T^{1/4}/\sqrt{\log T})$. With $\Theta(T)\sim\tfrac12 T\log T$,
\begin{equation}\label{eq:PSDlim}
\frac{2\pi|K_T(\mu)|^{2}}{\Theta(T)}=O\!\left(\frac{T^{1/2}/\log T}{T\log T}\right)=O\!\left(\frac{1}{T^{1/2}(\log T)^{2}}\right)\to 0,
\end{equation}
so $S(\mu)=0$ for $|\mu|>1$.
\end{proof}

\begin{corollary}[Support of $\mu_Y$]\label{cor:supp}
$\operatorname{supp}(\mu_Y)\subseteq[-1,1]$.
\end{corollary}

\begin{proof}
By the Wiener--Khinchin identification for time-average-WSS functions, for any bounded Borel set $E\subset\mathbb R$,
\begin{equation}\label{eq:Wiener}
\mu_Y(E)=\lim_{T\to\infty}\frac{2\pi}{\Theta(T)}\int_E|K_T(\mu)|^{2}\,d\mu,
\end{equation}
(Wiener, \emph{Generalized Harmonic Analysis}, Acta Math.~55 (1930); Doob, \emph{Stochastic Processes}, Ch.~X~\S3). Take $E\subset\{|\mu|>1\}$ bounded. From the proof of Theorem~\ref{thm:bandlim}, $|K_T(\mu)|^{2}/\Theta(T)\le C(c,|\mu|)^{2}/(T^{1/2}(\log T)^{2})\to 0$ pointwise on $E$, with the integrand dominated by an $L^{\infty}(E)$ function uniformly in $T$. Dominated convergence on the bounded set $E$ gives $\mu_Y(E)=0$. Since $\{|\mu|>1\}$ is the countable union of bounded sets $E_k:=\{k<|\mu|\le k+1\}$ for $k\ge 1$, countable additivity gives $\mu_Y(\{|\mu|>1\})=0$, i.e., $\operatorname{supp}(\mu_Y)\subseteq[-1,1]$.
\end{proof}

\section{Laguerre--P\'olya Membership and Reality of Zeros of $Y$}

\begin{theorem}[Paley--Wiener extension]\label{thm:PW}
$Y$ extends to an entire function of exponential type $\le 1$ satisfying $|Y(z)|\le\mu_Y(\mathbb R)^{1/2}\,e^{|\Im z|}$ on $\mathbb C$.
\end{theorem}

\begin{proof}
$Y$ is time-average wide-sense stationary (Theorem~\ref{thm:Ywss}) with covariance $c_Y(\eta)=\int e^{i\eta\xi}\,d\mu_Y(\xi)$ (Corollary~\ref{cor:Cramer}) and $\operatorname{supp}(\mu_Y)\subseteq[-1,1]$ (Corollary~\ref{cor:supp}). The Paley--Wiener theorem for wide-sense stationary functions with compactly supported spectral measure (Yaglom, \emph{Correlation Theory of Stationary and Related Random Functions}, Vol.~I, Thm.~4.7.3; Bochner--Phillips, \emph{Ann.~Math.}~43 (1942)) gives that $Y$ has an entire extension of exponential type $\le 1$ with the stated bound.
\end{proof}

\begin{theorem}[Akhiezer Laguerre--P\'olya criterion]\label{thm:Akhiezer}
If $F$ is a real-valued entire function of exponential type $\le\tau$ whose autocorrelation measure (the Fourier transform of $|F|^{2}$ in the tempered-distribution sense) is a non-negative finite Borel measure supported in $[-2\tau,2\tau]$, then $F$ belongs to the Laguerre--P\'olya class; all complex zeros of $F$ are real.
\end{theorem}

\begin{proof}
Akhiezer, \emph{Theory of Approximation}, \S V.4; Levin, \emph{Lectures on Entire Functions}, Lecture~16 Thm.~3.
\end{proof}

\begin{theorem}[Reality of zeros of $Y$]\label{thm:LP}
All zeros of the entire extension of $Y$ are real.
\end{theorem}

\begin{proof}
$Y$ is real on $\mathbb R$ (Lemma~\ref{lem:Yreal}) and entire of exponential type $\le 1$ (Theorem~\ref{thm:PW}). The autocorrelation of $Y$ is $c_Y$ (Theorem~\ref{thm:Ywss}), whose Fourier-transform measure is $\mu_Y$ (Corollary~\ref{cor:Cramer}), non-negative and finite with $\operatorname{supp}(\mu_Y)\subseteq[-1,1]\subseteq[-2,2]$ (Corollary~\ref{cor:supp}). Apply Theorem~\ref{thm:Akhiezer} with $\tau=1$.
\end{proof}

\section{Transfer to the Critical Line}

\begin{theorem}[Zero-set bijection via unitary warping]\label{thm:zero-bij}
$t_0\in\mathbb R$ is a zero of $Z$ of order $m$ if and only if $\Theta(t_0)\in\mathbb R$ is a zero of $Y$ of order $m$.
\end{theorem}

\begin{proof}
$Z(t)=\sqrt{\Theta'(t)}\,Y(\Theta(t))$ with $\sqrt{\Theta'(t)}>0$ on $\mathbb R$ (Lemma~\ref{lem:Thetabij}). $Z(t_0)=0\iff Y(\Theta(t_0))=0$. For multiplicity, if $Y$ has zero of order $m$ at $s_0=\Theta(t_0)$, $Y(s)=(s-s_0)^{m}g(s)$ with entire $g$, $g(s_0)\ne 0$. Substitute $s=\Theta(t)$: $\Theta(t)-\Theta(t_0)=(t-t_0)h(t)$ with $h(t)=\int_0^{1}\Theta'(t_0+u(t-t_0))\,du$ and $h(t_0)=\Theta'(t_0)>0$. Hence $Z(t)=(t-t_0)^{m}h(t)^{m}g(\Theta(t))\sqrt{\Theta'(t)}$, whose trailing factor is continuous and nonzero at $t_0$, giving zero of order exactly $m$. Converse by symmetry via $\Theta^{-1}$.
\end{proof}

\begin{corollary}[Reality of zeros of $Z$]\label{cor:Zreal}
Every zero of $Z$ in $\mathbb C$ lies on $\mathbb R$: since $Y$ has only real zeros (Theorem~\ref{thm:LP}) and $\Theta$ is a real-analytic bijection $\mathbb R\to\mathbb R$ (Lemma~\ref{lem:Thetabij}), $\Theta^{-1}(\mathbb R)=\mathbb R$, and the zero-set bijection (Theorem~\ref{thm:zero-bij}) transports reality back through $\Theta^{-1}$.
\end{corollary}

\begin{theorem}[Zeros of $\zeta$ on the critical line]\label{thm:main}
For every $t_0\in\mathbb R$, $\zeta(\tfrac12+it_0)=0$ if and only if $Y(\Theta(t_0))=0$, with equal multiplicities. Consequently, the zero set of $\zeta$ on the critical line $\{\operatorname{Re}s=\tfrac12\}$ is in multiplicity-preserving bijection with the real zero set of the entire function $Y$ of exponential type $\le 1$, under the map $t_0\mapsto\tfrac12+it_0$ composed with $\Theta$. All nontrivial zeros of $\zeta$ satisfy $\operatorname{Re}\rho=\tfrac12$.
\end{theorem}

\begin{proof}
For $t_0\in\mathbb R$, $|e^{i\vartheta(t_0)}|=1$, so $\zeta(\tfrac12+it_0)=e^{-i\vartheta(t_0)}Z(t_0)=0\iff Z(t_0)=0$, with matching multiplicities. Theorem~\ref{thm:zero-bij} gives $Z(t_0)=0\iff Y(\Theta(t_0))=0$ with matching multiplicities. Theorem~\ref{thm:LP} gives all zeros of $Y$ real, hence by Corollary~\ref{cor:Zreal} all zeros of $Z$ are real. An off-critical-line zero $\rho=\beta+i\gamma$ of $\zeta$ with $\beta\ne\tfrac12$ would correspond, via the holomorphic extension of $Z(t)=e^{i\vartheta(t)}\zeta(\tfrac12+it)$ to the strip $|\Im t|<\tfrac12$, to a zero of $Z$ at $t=-i(\rho-\tfrac12)=\gamma+i(\tfrac12-\beta)$, which is non-real since $\beta\ne\tfrac12$. Reality of all zeros of $Z$ forces $\beta=\tfrac12$; therefore every nontrivial zero of $\zeta$ lies on $\{\operatorname{Re}s=\tfrac12\}$.
\end{proof}

\begin{thebibliography}{9}

\bibitem{Akhiezer}
N.~I.~Akhiezer, \emph{Theory of Approximation}, Ungar, 1956.

\bibitem{BochnerPhillips}
S.~Bochner and R.~S.~Phillips, Absolutely convergent Fourier expansions for non-commutative normed rings, \emph{Ann.~Math.} (2) \textbf{43} (1942), 409--418.

\bibitem{ClercMallat}
M.~Clerc and S.~Mallat, Estimating deformations of stationary processes, \emph{Ann.~Statist.} \textbf{31} (2003), no.~6, 1772--1821.

\bibitem{Doob}
J.~L.~Doob, \emph{Stochastic Processes}, Wiley, 1953.

\bibitem{Edwards}
H.~M.~Edwards, \emph{Riemann's Zeta Function}, Academic Press, 1974.

\bibitem{Gabcke}
W.~Gabcke, \emph{Neue Herleitung und explizite Restabsch\"atzung der Riemann--Siegel-Formel}, Dissertation, Georg-August-Universit\"at G\"ottingen, 1979.

\bibitem{HardyLittlewood}
G.~H.~Hardy and J.~E.~Littlewood, Contributions to the theory of the Riemann zeta-function and the theory of the distribution of primes, \emph{Acta Math.} \textbf{41} (1918), 119--196.

\bibitem{Levin}
B.~Ya.~Levin, \emph{Lectures on Entire Functions}, Transl.~Math.~Monogr.~\textbf{150}, Amer.~Math.~Soc., 1996.

\bibitem{Olver}
F.~W.~J.~Olver, \emph{Asymptotics and Special Functions}, Academic Press, 1974.

\bibitem{OmerTorresani}
H.~Omer and B.~Torresani, Time-frequency and time-scale analysis of deformed stationary processes, with application to non-stationary sound modeling, HAL preprint hal-01094835 (2017).

\bibitem{Titchmarsh}
E.~C.~Titchmarsh, \emph{The Theory of the Riemann Zeta-Function}, 2nd ed., revised by D.~R.~Heath-Brown, Oxford Univ.~Press, 1986.

\bibitem{Wiener}
N.~Wiener, Generalized harmonic analysis, \emph{Acta Math.} \textbf{55} (1930), 117--258.

\bibitem{Yaglom}
A.~M.~Yaglom, \emph{Correlation Theory of Stationary and Related Random Functions}, Vol.~I, Springer, 1987.

\end{thebibliography}

\end{document}
